\chapter*{TÓM TẮT}
\vspace{-40pt}

Đề tài ``Xây dựng website thương mại điện tử bán máy tính và linh kiện máy tính'' tập trung xây dựng kênh trực tuyến để cửa hàng T\&T Computer giới thiệu sản phẩm và hỗ trợ khách hàng tra cứu, đặt mua máy tính cùng linh kiện. Hệ thống hướng đến khách hàng và các nhân sự quản trị hoạt động bán hàng. Website được xây dựng bằng PHP, sử dụng MySQL và PDO để lưu trữ, truy xuất dữ liệu; giao diện sử dụng HTML, CSS và JavaScript.

Các chức năng được tổ chức theo những nhóm nghiệp vụ chính gồm tra cứu sản phẩm; đăng ký, đăng nhập và quản lý tài khoản; quản lý giỏ hàng, địa chỉ nhận hàng và đơn hàng; cùng các chức năng dành cho nhân viên và quản trị viên. Khu vực quản trị hỗ trợ quản lý sản phẩm, danh mục, người dùng, tồn kho, khuyến mãi, voucher, đơn hàng và thống kê. Quy trình đơn hàng bao gồm kiểm tra tồn kho, áp dụng ưu đãi, theo dõi trạng thái và xử lý hủy hoặc giao hàng thất bại theo các quy tắc được mô tả trong báo cáo.

Về kỹ thuật, ứng dụng tổ chức xử lý qua front controller và router, với các thành phần controller, model và view. Báo cáo trình bày sơ đồ Use Case, sơ đồ trình tự, sơ đồ lớp và thiết kế cơ sở dữ liệu để mô tả yêu cầu và cấu trúc hệ thống. Mã nguồn có hỗ trợ đăng nhập và liên kết tài khoản Google; biểu mẫu đặt hàng có lựa chọn COD nhưng chưa có thao tác ghi nhận thu tiền, còn thanh toán trực tuyến được tích hợp với SePay Sandbox qua checkout và IPN. Tích hợp SePay chỉ ở môi trường thử nghiệm, chưa được xem là vận hành thanh toán thực tế.

Việc đánh giá trong báo cáo chủ yếu dựa trên đối chiếu mã nguồn, route, controller, model, giao diện và các luồng xử lý. Báo cáo không nêu số lượng ca kiểm thử hoặc tỷ lệ đạt; đồng thời chưa ghi nhận kiểm thử trực tiếp trên Apache và MySQL với dữ liệu thực. Vì vậy, kết quả hiện tại phản ánh mức độ hoàn thiện chức năng và các tình huống đã xem xét trong phạm vi đồ án, chưa thay thế cho kiểm thử hệ thống thực tế hoặc đánh giá vận hành dài hạn.

Nhóm đã xây dựng được các luồng mua hàng và quản trị cơ bản, đồng thời vận dụng kiến thức phân tích, thiết kế cơ sở dữ liệu và phát triển ứng dụng web vào bài toán cụ thể. Hệ thống còn cần được kiểm thử trên môi trường triển khai, hoàn thiện cấu hình và giám sát thanh toán, cũng như tiếp tục cải thiện giao diện và khả năng mở rộng trước khi xem xét áp dụng trong thực tế.

